\documentclass[11pt]{article}
\usepackage{amsmath,amssymb,amsthm,mathtools}
\usepackage[margin=1in]{geometry}
\usepackage{hyperref}

\newtheorem{theorem}{Theorem}
\newtheorem{proposition}[theorem]{Proposition}
\newtheorem{lemma}[theorem]{Lemma}
\newtheorem{corollary}[theorem]{Corollary}
\theoremstyle{definition}
\newtheorem{remark}[theorem]{Remark}
\newtheorem{hypothesis}[theorem]{Hypothesis}

\newcommand{\bd}{\beta_d}
\newcommand{\bdm}{\beta_{d-1}}
\newcommand{\rinf}{r_\infty}
\newcommand{\N}{\mathbb{N}}
\newcommand{\C}{\mathbb{C}}
\newcommand{\R}{\mathbb{R}}

\title{The radius and leading growth of polynomial continued-fraction
numerators\\[4pt]
\large Session \texttt{CT-S1-THEOREM}, prompt P4 Task A}
\author{}
\date{}

\begin{document}
\maketitle

\begin{abstract}
Let $b\in\C[n]$ have degree $d\ge 1$ and let $Q_n=b(n)Q_{n-1}+Q_{n-2}$ with
$Q_{-1}=0$, $Q_0=1$. We show that $r_n:=Q_n/\prod_{k\le n}b(k)$ always
converges (Proposition~\ref{prop:exists}), that its limit $\rinf$ is nonzero
whenever $b(n)>0$ for all $n\ge1$ (Proposition~\ref{prop:positive}), and that
under the hypothesis $\rinf\neq0$ the sequence $Q_n$ has radius
$\lim|Q_n/(dn)!|^{1/n}=|\bd|/d^d$ (Theorem~\ref{thm:radius}) and leading growth
$Q_n/\Gamma(dn+1)=K(\bd/d^d)^n n^s(1+O(1/n))$ with
$s=\bdm/\bd+(d-1)/2$ and an explicit $K$ (Theorem~\ref{thm:growth}).
Every step is marked \textsc{complete} or \textsc{gap}.
\textbf{Overall label: Theorems~\ref{thm:radius} and~\ref{thm:growth} are both
\textsc{proven}} --- Theorem~\ref{thm:radius} with no imported input at all, and
Theorem~\ref{thm:growth} modulo the single classical Gamma-ratio normalisation
DLMF~(5.11.12), which is named explicitly at the point of
use. (In the first draft Theorem~\ref{thm:growth} was \textsc{structural}; Gap~G1
was closed in audit round R1, and a load-bearing unproved assertion in
Theorem~\ref{thm:radius}, Step~3, was supplied with a proof at the same time.)
\end{abstract}

\section{Setting}

\begin{equation}\label{eq:setting}
b(n)=\sum_{k=0}^{d}\beta_k n^k=\bd\prod_{i=1}^{d}(n-r_i),
\qquad \bd\neq0,\quad d\ge1,
\end{equation}
with $r_1,\dots,r_d\in\C$ the roots listed with multiplicity. Standing
assumptions:
\begin{itemize}
\item[\textbf{(A1)}] $b(k)\neq0$ for every integer $k\ge1$.
\item[\textbf{(A2)}] No $r_i$ is a positive integer, equivalently
$\Gamma(1-r_i)$ is finite for every $i$. (Recall $\Gamma$ has poles exactly at
$0,-1,-2,\dots$.) \emph{(A2) is used only in Theorem~\ref{thm:growth}; note
(A2) $\Rightarrow$ (A1).}
\end{itemize}
Define
\begin{equation}\label{eq:rec}
Q_{-1}=0,\qquad Q_0=1,\qquad Q_n=b(n)Q_{n-1}+Q_{n-2}\ \ (n\ge1),
\end{equation}
and, using (A1),
\begin{equation}\label{eq:rn}
B_n:=\prod_{k=1}^{n}b(k),\qquad r_n:=\frac{Q_n}{B_n}\quad(n\ge0),\qquad B_0:=1 .
\end{equation}

\section{The hypothesis (H) and its status}

\begin{proposition}[Existence of the limit; unconditional]\label{prop:exists}
Under \textup{(A1)} the limit $\rinf:=\lim_{n\to\infty}r_n$ exists in $\C$.
Moreover
\begin{equation}\label{eq:rate}
\rinf-r_n=O\!\left(n^{-(2d-1)}\right).
\end{equation}
\end{proposition}

\begin{proof}
Divide \eqref{eq:rec} by $B_n=b(n)B_{n-1}=b(n)b(n-1)B_{n-2}$:
\begin{equation}\label{eq:rnrec}
r_n=r_{n-1}+\frac{r_{n-2}}{b(n-1)b(n)}\qquad(n\ge2),
\end{equation}
with $r_0=1$ and $r_1=Q_1/b(1)=b(1)/b(1)=1$ (using $Q_{-1}=0$).
\textsc{complete.}

Since $\deg b=d$ with leading coefficient $\bd\neq0$, there are constants
$c>0$ and $n_0$ with $|b(n)|\ge c\,n^{d}$ for $n\ge n_0$; hence for $n\ge n_0+1$
\begin{equation}\label{eq:eps}
\varepsilon_n:=\frac{1}{|b(n-1)b(n)|}\le \frac{C}{n^{2d}},\qquad C:=2^{d}c^{-2},
\end{equation}
where the factor $2^d$ absorbs $(n-1)^{-d}\le 2^d n^{-d}$ for $n\ge2$.
As $2d\ge2$, $\sum_n\varepsilon_n<\infty$. \textsc{complete.}

Let $M_n:=\max(|r_n|,|r_{n-1}|)$. From \eqref{eq:rnrec},
$|r_n|\le|r_{n-1}|+\varepsilon_n|r_{n-2}|\le(1+\varepsilon_n)M_{n-1}$, so
$M_n\le(1+\varepsilon_n)M_{n-1}$ and therefore
\[
M_n\le M_{n_0}\prod_{k>n_0}(1+\varepsilon_k)\le
M_{n_0}\exp\Big(\sum_{k>n_0}\varepsilon_k\Big)=:M<\infty .
\]
So $(r_n)$ is bounded. \textsc{complete.}

Consequently $|r_n-r_{n-1}|=\varepsilon_n|r_{n-2}|\le M\varepsilon_n$, and
$\sum_n|r_n-r_{n-1}|\le M\sum_n\varepsilon_n<\infty$: the telescoping series
$r_0+\sum_{n\ge1}(r_n-r_{n-1})$ converges absolutely, so $r_n$ converges to some
$\rinf\in\C$. \textsc{complete.}

For \eqref{eq:rate}, $|\rinf-r_n|\le\sum_{k>n}|r_k-r_{k-1}|
\le MC\sum_{k>n}k^{-2d}=O(n^{-(2d-1)})$, the last step by comparison with
$\int_n^\infty x^{-2d}\,dx=n^{-(2d-1)}/(2d-1)$, valid since $2d\ge2>1$.
\textsc{complete.}
\end{proof}

\begin{hypothesis}[H]\label{hyp:H}
$\rinf\neq0$.
\end{hypothesis}

\begin{proposition}[(H) is automatic in the positive case]\label{prop:positive}
If $b(n)>0$ for every integer $n\ge1$, then $r_n$ is a strictly increasing
sequence of positive reals with $r_n\ge1$, hence $\rinf\ge1>0$ and
\textup{(H)} holds.
\end{proposition}

\begin{proof}
By induction $r_n>0$ for all $n\ge0$: $r_0=r_1=1>0$, and if
$r_{n-1},r_{n-2}>0$ then \eqref{eq:rnrec} gives
$r_n=r_{n-1}+r_{n-2}/(b(n-1)b(n))>r_{n-1}>0$ since $b(n-1)b(n)>0$.
Thus $(r_n)_{n\ge1}$ is strictly increasing with $r_1=1$, so $r_n\ge1$ for
$n\ge1$; combined with Proposition~\ref{prop:exists} the limit exists and
satisfies $\rinf\ge1$. \textsc{complete.}
\end{proof}

\begin{remark}[Pincherle interpretation of (H)]\label{rem:pincherle}
Pincherle's theorem (standard; see e.g.\ Gautschi, \emph{Computational aspects
of three-term recurrence relations}, SIAM Rev.\ 9 (1967), and Lorentzen--Waadeland,
\emph{Continued Fractions with Applications}, Ch.~IV) states that the continued
fraction associated with a three-term recurrence converges if and only if the
recurrence possesses a \emph{minimal} solution, i.e.\ a solution $y_n$ with
$y_n/z_n\to0$ for every linearly independent solution $z_n$.

Here $\prod_{k\le n}b(k)$ is, up to a constant, the natural ``dominant'' scale.
Hypothesis (H) says precisely that $Q_n$ is \emph{not} the minimal solution of
\eqref{eq:rec}: the minimal solution $P_n$ satisfies $P_n/B_n\to0$, so a
solution with $Q_n/B_n\to\rinf\neq0$ is dominant. Equivalently, (H) fails exactly
when the chosen initial data $(Q_{-1},Q_0)=(0,1)$ happens to land on the
one-dimensional minimal subspace. Since that subspace has codimension one, (H) is
generic; Proposition~\ref{prop:positive} makes it automatic in the positive case.
\textbf{Label: \textsc{structural}} --- Pincherle's theorem is cited as standard
and the identification of $B_n$ with the dominant scale is asserted here only for
interpretation; nothing below depends on this remark.
\end{remark}

\section{Part (i): the radius}

\begin{theorem}[Radius]\label{thm:radius}
Assume \textup{(A1)} and \textup{(H)}. Then
\begin{equation}\label{eq:radius}
\lim_{n\to\infty}\left|\frac{Q_n}{(dn)!}\right|^{1/n}=\frac{|\bd|}{d^{\,d}} .
\end{equation}
Consequently the Borel-type transform
\begin{equation}\label{eq:borel}
\mathcal{B}(w):=\sum_{n\ge0}\frac{Q_n}{(dn)!}\,w^n
\qquad\bigl((dn)!=\Gamma(dn+1)\bigr)
\end{equation}
has radius of convergence $d^{\,d}/|\bd|$ in the variable $w$, and in the
uniformiser $\zeta=w^{1/d}$ the corresponding radius is
\begin{equation}\label{eq:xi0}
\xi_0=\frac{d}{|\bd|^{1/d}} .
\end{equation}
\end{theorem}

\begin{remark}[Normalisation, stated explicitly]\label{rem:norm}
Three conventions are fixed once and for all and must be quoted with any use of
\eqref{eq:xi0}:
\begin{enumerate}
\item the \emph{series} transformed is $\sum_n Q_n w^n/\Gamma(dn+1)$, i.e.\ the
coefficients are the recurrence numerators $Q_n$ themselves, divided by
$\Gamma(dn+1)$ --- \emph{not} by $n!$, $\Gamma(n+1)$, $\Gamma(2n+1)$ or
$\Gamma(2n)$;
\item the \emph{Borel variable} is $w$, conjugate to the index $n$;
\item the \emph{uniformiser} is $\zeta=w^{1/d}$, so that $\xi_0=|w^*|^{1/d}$
with $w^*=d^d/\bd$.
\end{enumerate}
At $d=2$ this gives $\xi_0=2/\sqrt{\beta_2}$, matching the statement labelled
\texttt{prop:xi0} in Channel Theory v1.3.1.
\end{remark}

\begin{proof}[Proof of Theorem~\ref{thm:radius}]
We deliberately avoid Gamma asymptotics, using only $r_n\to\rinf\neq0$ and
elementary bounds, together with $\lim_{n}(m!)^{1/n}$-type facts in the
elementary form below.

\medskip\noindent\textbf{Step 1 (reduce to $B_n$).}
By \eqref{eq:rn}, $Q_n=r_nB_n$ exactly. Since $r_n\to\rinf\neq0$, we have
$|r_n|^{1/n}\to1$: indeed for large $n$, $\tfrac12|\rinf|\le|r_n|\le2|\rinf|$,
and $c^{1/n}\to1$ for any fixed $c\in(0,\infty)$. Hence
$\lim|Q_n/(dn)!|^{1/n}=\lim|B_n/(dn)!|^{1/n}$ whenever the latter exists.
\emph{This is the only place (H) is used, and it is used essentially: if
$\rinf=0$ then $|r_n|^{1/n}$ need not tend to $1$.} \textsc{complete.}

\medskip\noindent\textbf{Step 2 ($B_n$ against $(n!)^d$).}
Write $b(k)=\bd k^d\,u_k$ with $u_k:=b(k)/(\bd k^d)\to1$. Then
\[
B_n=\bd^{\,n}\,(n!)^{d}\prod_{k=1}^{n}u_k .
\]
Since $u_k=1+O(1/k)$ and $u_k\neq0$ by (A1), we have $\log|u_k|=O(1/k)$, so
$\bigl|\sum_{k\le n}\log|u_k|\bigr|=O(\log n)$ and therefore
$\bigl(\prod_{k\le n}|u_k|\bigr)^{1/n}=\exp\bigl(O(\log n)/n\bigr)\to1$.
Hence
\begin{equation}\label{eq:Bnroot}
\lim_{n\to\infty}\left|\frac{B_n}{(n!)^{d}}\right|^{1/n}=|\bd| .
\end{equation}
\textsc{complete.} (Only $\log(1+x)=O(x)$ and $\sum_{k\le n}1/k=O(\log n)$ are
used; no Stirling.)

\medskip\noindent\textbf{Step 3 ($(n!)^d$ against $(dn)!$).}
We claim
\begin{equation}\label{eq:multinomial}
\lim_{n\to\infty}\left(\frac{(dn)!}{(n!)^{d}}\right)^{1/n}=d^{\,d}.
\end{equation}
Write $C(k_1,\dots,k_d):=\binom{dn}{k_1,\dots,k_d}=(dn)!/(k_1!\cdots k_d!)$ for
a composition $k_1+\cdots+k_d=dn$ into non-negative parts.

\emph{Sub-lemma (the balanced composition is maximal).} $C$ attains its maximum
over all such compositions at $k_1=\cdots=k_d=n$. Indeed, suppose
$k_i\ge k_j+2$ for some $i\neq j$. Replacing $(k_i,k_j)$ by $(k_i-1,k_j+1)$
leaves the total $dn$ unchanged and multiplies $C$ by
\[
\frac{k_i!\,k_j!}{(k_i-1)!\,(k_j+1)!}=\frac{k_i}{k_j+1}\ \ge\ \frac{k_j+2}{k_j+1}>1 ,
\]
so $C$ strictly increases. As the set of compositions is finite, any maximiser
must therefore satisfy $|k_i-k_j|\le1$ for all $i,j$; together with
$\sum_ik_i=dn$ this forces $k_1=\cdots=k_d=n$. Hence
$(dn)!/(n!)^d=\max C$. \hfill$\square$

Now $(dn)!/(n!)^d$ is one term of the multinomial expansion
$\sum_{k_1+\cdots+k_d=dn}C(k_1,\dots,k_d)=(\underbrace{1+\cdots+1}_{d})^{dn}=d^{\,dn}$,
which gives the upper bound; and being the \emph{maximum} it is at least the
average, i.e.\ $d^{\,dn}$ divided by the number $N_{d,n}$ of compositions. Each
composition is determined by $(k_1,\dots,k_{d-1})\in\{0,1,\dots,dn\}^{d-1}$
(the last part being forced), so $N_{d,n}\le(dn+1)^{d-1}$. Therefore
\[
\frac{d^{\,dn}}{(dn+1)^{d-1}}\ \le\ \frac{(dn)!}{(n!)^d}\ \le\ d^{\,dn}.
\]
Taking $n$-th roots and using $(dn+1)^{(d-1)/n}\to1$ gives
\eqref{eq:multinomial}. \textsc{complete.} (Purely combinatorial; no Stirling,
as required.)

\medskip\noindent\textbf{Step 4 (combine).}
Multiplying the limits from Steps 1--3,
\[
\lim_n\left|\frac{Q_n}{(dn)!}\right|^{1/n}
=\lim_n\left|\frac{B_n}{(n!)^{d}}\right|^{1/n}
\Big/\lim_n\left(\frac{(dn)!}{(n!)^{d}}\right)^{1/n}
=\frac{|\bd|}{d^{\,d}},
\]
which is \eqref{eq:radius}. \textsc{complete.}

\medskip\noindent\textbf{Step 5 (radius and uniformiser).}
By Cauchy--Hadamard the radius of \eqref{eq:borel} is the reciprocal of
\eqref{eq:radius}, namely $d^d/|\bd|$. Under $\zeta=w^{1/d}$ the disc
$|w|<d^d/|\bd|$ is $|\zeta|<\bigl(d^d/|\bd|\bigr)^{1/d}=d/|\bd|^{1/d}=\xi_0$.
\textsc{complete.}
\end{proof}

\begin{corollary}[$d=2$, V\_quad]\label{cor:vquad}
For $b(n)=3n^2+n+1$ one gets $\xi_0=2/\sqrt3$, in agreement with the stored
value in \texttt{sessions/2026-05-01/CC-PIPELINE-G/results\_vquad.json}.
\end{corollary}

\section{Part (ii): the leading growth}

\begin{theorem}[Growth]\label{thm:growth}
Assume \textup{(A2)} and \textup{(H)}. Then, with
\begin{equation}\label{eq:sK}
s=\frac{\bdm}{\bd}+\frac{d-1}{2},
\qquad
K=\rinf\cdot\frac{(2\pi)^{(d-1)/2}}{\sqrt{d}\ \prod_{i=1}^{d}\Gamma(1-r_i)},
\end{equation}
one has
\begin{equation}\label{eq:growth}
\frac{Q_n}{\Gamma(dn+1)}
=K\left(\frac{\bd}{d^{\,d}}\right)^{\!n} n^{s}\bigl(1+O(1/n)\bigr).
\end{equation}
\end{theorem}

\begin{proof}
\textbf{Step 1 (exact product formula).} By (A2) no $r_i$ is a positive
integer, so telescoping $\Gamma(x+1)=x\Gamma(x)$ gives, for each $i$,
$\prod_{k=1}^{n}(k-r_i)=\Gamma(n+1-r_i)/\Gamma(1-r_i)$ with both Gamma values
finite and nonzero. Hence, by \eqref{eq:setting},
\begin{equation}\label{eq:Bnexact}
B_n=\bd^{\,n}\prod_{i=1}^{d}\frac{\Gamma(n+1-r_i)}{\Gamma(1-r_i)} .
\end{equation}
\textsc{complete} (exact, elementary).

\medskip\noindent\textbf{Step 2 (convergence of the correction product).}
This is the step the prompt asks to be exhibited. Put $\kappa:=\bdm/\bd$ and
consider
\[
\Pi_n:=\prod_{k=2}^{n}\left[\frac{b(k)}{\bd k^{d}}\left(\frac{k}{k-1}\right)^{-\kappa}\right].
\]
Now $b(k)/(\bd k^d)=\prod_i(1-r_i/k)=1-\bigl(\sum_i r_i\bigr)/k+O(1/k^2)
=1+\kappa/k+O(1/k^2)$, using Vieta in the form
\begin{equation}\label{eq:vieta}
\sum_{i=1}^{d}r_i=-\frac{\bdm}{\bd}=-\kappa .
\end{equation}
Also $\bigl(k/(k-1)\bigr)^{-\kappa}=(1-1/k)^{\kappa}=1-\kappa/k+O(1/k^2)$.
Multiplying, the bracket equals
$\bigl(1+\kappa/k+O(1/k^2)\bigr)\bigl(1-\kappa/k+O(1/k^2)\bigr)=1+O(1/k^2)$.
A product $\prod(1+a_k)$ with $\sum|a_k|<\infty$ converges to a nonzero limit;
here $\sum_k O(1/k^2)<\infty$. Hence $\Pi_n$ converges to a finite nonzero
limit. \textsc{complete.}

This is exactly the mechanism that produces the exponent $\kappa=\bdm/\bd$:
the telescoping factor $\prod_{k\le n}(k/(k-1))^{\kappa}=n^{\kappa}$ carries it.

\medskip\noindent\textbf{Step 3 (Gamma-ratio asymptotics).}
For fixed $r\in\C$,
\begin{equation}\label{eq:gammaratio}
\frac{\Gamma(n+1-r)}{\Gamma(n+1)}=n^{-r}\bigl(1+O(1/n)\bigr).
\end{equation}

\emph{Proof of \eqref{eq:gammaratio}.} Put
$h(n):=\Gamma(n+1-r)\,n^{r}/\Gamma(n+1)$ for $n>\max(0,\Re r)$, so that
\eqref{eq:gammaratio} is the assertion $h(n)=1+O(1/n)$. Using
$\Gamma(x+1)=x\Gamma(x)$ twice,
\begin{equation}\label{eq:hratio}
\frac{h(n)}{h(n-1)}
=\frac{\Gamma(n+1-r)}{\Gamma(n-r)}\cdot\frac{\Gamma(n)}{\Gamma(n+1)}
\cdot\Bigl(\frac{n}{n-1}\Bigr)^{r}
=\Bigl(1-\frac{r}{n}\Bigr)\Bigl(1-\frac1n\Bigr)^{-r} .
\end{equation}
Expanding, $(1-1/n)^{-r}=1+r/n+O(1/n^{2})$, so the right-hand side of
\eqref{eq:hratio} equals $(1-r/n)(1+r/n+O(1/n^2))=1+O(1/n^{2})$. Since
$\sum_n O(1/n^2)<\infty$, the telescoping product
$h(n)=h(n_0)\prod_{m=n_0+1}^{n}h(m)/h(m-1)$ converges to a finite nonzero limit
$h_\infty$, and because the tail of a convergent $\sum O(1/m^{2})$ is $O(1/n)$,
\begin{equation}\label{eq:hlimit}
h(n)=h_\infty\bigl(1+O(1/n)\bigr).
\end{equation}
It remains to identify $h_\infty=1$. This is the classical Gamma-ratio
normalisation: DLMF~(5.11.12) states
$\Gamma(z+a)/\Gamma(z+b)\sim z^{a-b}$ as
$|z|\to\infty$ with $|\operatorname{ph}z|\le\pi-\delta$
\cite[\S5.11(iii)]{DLMF}. Taking $a=1-r$, $b=1$, $z=n$ gives
$\Gamma(n+1-r)/\Gamma(n+1)\sim n^{-r}$, i.e.\ $h_\infty=1$. (The refined
expansion $\Gamma(z+a)/\Gamma(z+b)\sim z^{a-b}\sum_{k\ge0}G_k(a,b)z^{-k}$ is
DLMF~(5.11.13), with $G_0(a,b)=1$ at DLMF~(5.11.15); only the leading statement
(5.11.12) is needed here.) Substituting into
\eqref{eq:hlimit} yields \eqref{eq:gammaratio}. \hfill$\square$

Applying \eqref{eq:gammaratio} for each $i$ and multiplying, then using
\eqref{eq:vieta},
\begin{equation}\label{eq:prodgamma}
\prod_{i=1}^{d}\Gamma(n+1-r_i)=\Gamma(n+1)^{d}\,n^{\kappa}\bigl(1+O(1/n)\bigr).
\end{equation}
\textsc{complete}, modulo the single cited classical input $h_\infty=1$; see
\S\ref{sec:gaps}. (Step 2 is an independent, elementary route to the same
exponent $n^{\kappa}$, and \eqref{eq:hlimit} above supplies the $O(1/n)$
bookkeeping without any citation, so the cited input is only the
\emph{constant}.)

\medskip\noindent\textbf{Step 4 (Gauss multiplication).}
The Gauss multiplication formula
$\Gamma(dz)=(2\pi)^{(1-d)/2}d^{\,dz-1/2}\prod_{j=0}^{d-1}\Gamma(z+j/d)$ at
$z=n$, together with \eqref{eq:gammaratio} in the form
$\Gamma(n+j/d)=\Gamma(n)n^{j/d}(1+O(1/n))$ and
$\sum_{j=0}^{d-1}j/d=(d-1)/2$, gives
$\Gamma(dn)=(2\pi)^{(1-d)/2}d^{\,dn-1/2}\Gamma(n)^{d}n^{(d-1)/2}(1+O(1/n))$.
Multiplying by $dn$ and using $\Gamma(n+1)^d=n^d\Gamma(n)^d$,
\begin{equation}\label{eq:gauss}
\Gamma(dn+1)=(2\pi)^{(1-d)/2}d^{\,dn+1/2}\,\Gamma(n+1)^{d}\,n^{-(d-1)/2}
\bigl(1+O(1/n)\bigr).
\end{equation}
\textsc{complete}, modulo the same single cited constant recorded in
\S\ref{sec:gaps}.

\medskip\noindent\textbf{Step 5 (assemble).}
By Proposition~\ref{prop:exists} and \eqref{eq:rate}, $r_n=\rinf(1+O(1/n))$
(here (H) is used to divide by $\rinf$), so $Q_n=r_nB_n=\rinf B_n(1+O(1/n))$.
Combining with \eqref{eq:Bnexact} and \eqref{eq:prodgamma},
\[
Q_n=\rinf\,\bd^{\,n}\,
\frac{\Gamma(n+1)^{d}\,n^{\kappa}}{\prod_i\Gamma(1-r_i)}\bigl(1+O(1/n)\bigr).
\]
Dividing by \eqref{eq:gauss}, the factors $\Gamma(n+1)^{\pm d}$ cancel
\emph{exactly} --- this cancellation is the reason for normalising by
$\Gamma(dn+1)$ rather than by $n!$ --- leaving
\[
\frac{Q_n}{\Gamma(dn+1)}
=\frac{\rinf(2\pi)^{(d-1)/2}}{\sqrt d\prod_i\Gamma(1-r_i)}
\left(\frac{\bd}{d^{\,d}}\right)^{\!n}
n^{\kappa+(d-1)/2}\bigl(1+O(1/n)\bigr),
\]
which is \eqref{eq:growth} with \eqref{eq:sK}. \textsc{complete.}
\end{proof}

\begin{corollary}[Consistency with Theorem~\ref{thm:radius}]
\eqref{eq:growth} implies \eqref{eq:radius}, since $|K|^{1/n}\to1$ and
$n^{s/n}\to1$. The two parts are consistent, and part (i) is strictly weaker but
proved without Gap~G1.
\end{corollary}

\section{Status of each step, and the gaps}\label{sec:gaps}

\begin{itemize}
\item Proposition~\ref{prop:exists}: all steps \textsc{complete}.
\textbf{\textsc{proven}.}
\item Proposition~\ref{prop:positive}: \textsc{complete}. \textbf{\textsc{proven}.}
\item Remark~\ref{rem:pincherle}: cites Pincherle as standard; interpretive only.
\textbf{\textsc{structural}.}
\item Theorem~\ref{thm:radius}: Steps 1--5 all \textsc{complete}; no Gamma
asymptotics used, as required. \textbf{\textsc{proven}.}
\item Theorem~\ref{thm:growth}: Steps 1--5 all \textsc{complete}. The only
imported ingredient is the classical constant $h_\infty=1$ in
\eqref{eq:gammaratio}, cited to DLMF~(5.11.12); everything else, including the
$O(1/n)$ error bookkeeping, is proved here. \textbf{\textsc{proven}, modulo one
explicitly named classical citation.}
\end{itemize}

\paragraph{Gap G1 --- now closed (audit round R1).} In the first draft the whole
of \eqref{eq:gammaratio} was quoted rather than proved, and
Theorem~\ref{thm:growth} was accordingly labelled \textsc{structural}. Step~3
above now proves, from $\Gamma(x+1)=x\Gamma(x)$ alone, that
$h(n)=\Gamma(n+1-r)n^{r}/\Gamma(n+1)$ satisfies $h(n)/h(n-1)=1+O(1/n^{2})$ and
hence $h(n)=h_\infty(1+O(1/n))$ for a finite nonzero $h_\infty$. The residual
dependence is the single normalisation $h_\infty=1$, taken from
DLMF~(5.11.12). This is a citation of a classical result, not an
open gap; it is flagged here so that any reader can see exactly what is
imported. \textbf{No other gap remains.}

\paragraph{Audit note on Theorem~\ref{thm:radius}, Step 3.} The first draft
asserted without proof that the balanced multinomial coefficient is the largest
one. That assertion is load-bearing --- it is what converts ``at least the
average'' into the lower bound $d^{\,dn}/(dn+1)^{d-1}$ --- and it is now proved
by the exchange argument in the sub-lemma of Step~3.

\paragraph{Numerical corroboration (not a substitute for proof).}
$s$ in \eqref{eq:sK} was measured for eight degree-$4$ representatives in session
\texttt{CT-S2-RHO-D4} with worst residual $7.12\times10^{-16}$, and $K$ in
\eqref{eq:sK} was confirmed to $109.7$--$114.7$ digits over eight polynomials
spanning $d=2,3,4$ in session \texttt{CT-S3B-RINF}. Per the governance rules
these are \textsc{verified}, never \textsc{proven}.

\section{Scope remark}

Stated plainly, and deliberately restrictive:

\begin{enumerate}
\item \textbf{The theorem gives the radius and the leading growth only.} It
locates the singularity of $\mathcal{B}(w)$ nearest the origin and gives the
leading coefficient asymptotics of $Q_n/\Gamma(dn+1)$. It does \emph{not}
determine the \emph{analytic type} of the Borel singularity at $\xi_0$: whether
it is a pure branch point of exponent $-(s+1)$, whether logarithms occur, and
what the subleading germ looks like are all outside the statement.
\item \textbf{It reads only $\bd$ and $\bdm$.} The exponent $s$ depends on no
other coefficient, and the radius depends only on $\bd$. The lower coefficients
$\beta_0,\dots,\beta_{d-2}$ enter \eqref{eq:growth} \emph{only} through the
scalar prefactor $K$ (via $\rinf$ and $\prod_i\Gamma(1-r_i)$). Consequently any
test that inspects only the radius or only the exponent is insensitive to
$\beta_0,\dots,\beta_{d-2}$ and will return the same verdict on an entire family
of $b$'s. This is the structural explanation of a vacuity phenomenon of exactly
this shape.
\item \textbf{A pointer the operator must supply or correct.} The prompt asks
this scope remark to reference the open item \texttt{op:xi0-singularity-type}
and ``the v1.6 vacuity finding''. \emph{Neither could be located in this
repository} (see \texttt{sessions/2026-09-17/CT-S0-INVENTORY/REPORT.md}, \S3
items 4 and 5: the identifier \texttt{op:xi0-singularity-type} does not occur,
and the newest in-repo Channel Theory source is v1.3.1, with v1.5/v1.6 living in
a \texttt{pcf-research/} repository not present on this machine). Per the
governance rule against inventing identifiers, no citation is fabricated; item
(1) above states the scope limitation on its own terms, and item (2) states the
vacuity mechanism on its own terms. The operator should supply the correct
identifiers if these are to be cross-referenced.
\item \textbf{Relation to $\zeta^*=4/\sqrt3$.} Session
\texttt{CT-S3A-ZETA-AUDIT} concluded (classification (b)) that $\zeta^*=2\xi_0$
is a \emph{different} singularity --- the action difference $|c_+-c_-|$ of the
two formal solutions --- and not a re-expression of $\xi_0$ in another variable.
Nothing in this theorem bears on $\zeta^*$.
\end{enumerate}

\begin{thebibliography}{9}
\bibitem{DLMF}
NIST Digital Library of Mathematical Functions, \texttt{https://dlmf.nist.gov/},
Version 1.2.8; release date 2026-09-15 (accessed 2026-09-17).
\S5.11(iii) ``Ratios'', Eq.~(5.11.12):
$\Gamma(z+a)/\Gamma(z+b)\sim z^{a-b}$ for
$|\operatorname{ph}z|\le\pi-\delta$. (The expansion with coefficients
$G_k(a,b)$ is Eq.~(5.11.13), and $G_0(a,b)=1$ is Eq.~(5.11.15).)
\end{thebibliography}

\end{document}
